% concatenated from contractibilitiy-of-e-nets-in-Rn-ArXiv.tex
\documentclass{article}
\usepackage[utf8]{inputenc}
\usepackage[T2A]{fontenc}
\usepackage[english]{babel}


% \usepackage[a5paper,margin=0.8cm,bmargin=1.5cm]{geometry}
\usepackage{parskip}

\usepackage{amsmath}
\usepackage{amssymb}
\usepackage{amsthm}
\usepackage{amsfonts}
\usepackage{asymptote}
% \usepackage{hyperref}
\usepackage{xcolor}
\usepackage{subcaption}
\usepackage{graphicx}
\usepackage{cancel}


\newtheorem{problem}{}[]
%\let%\oldthm\thm
%\renewcommand{\thm}{\oldthm\normalfont}

\def\labelenumi{\rm(\theenumi)}

\textwidth=510pt
\textheight=670pt%
\oddsidemargin=15pt%
\evensidemargin=15pt
\topmargin=-11pt \headsep=14pt \headheight=22pt \voffset=-28pt
\hoffset=-30pt
\clubpenalty=10000
\widowpenalty=10000



% \newenvironment{answer}{\par\emph{Ответ:}}{\par}
% \newtheorem{lemma}{Лемма}[]
% \newtheorem{claim}{Утверждение}[]
% \newtheorem{remark}{Замечание}[]
% \newtheorem{conjecture}{Гипотеза}[]
% \newtheorem{proposition}{Предложение}[]
% \newtheorem{FunFact}{FunFact}[]
% \newtheorem{theorem}{Теорема}[]
% \newtheorem{corollary}{Следствие}[]
% \theoremstyle{definition}
% \newtheorem{definition}{Определение}[]
% \newtheorem{example}{Пример}[]
% \theoremstyle{plain}
\newenvironment{answer}{\par\emph{Ответ:}}{\par}
\newtheorem{lemma}{Lemma}[]
\newtheorem{claim}{Claim}[]
\newtheorem{remark}{Remark}[]
\newtheorem{conjecture}{Conjecture}[]
\newtheorem{proposition}{Proposition}[]
\newtheorem{FunFact}{FunFact}[]
\newtheorem{theorem}{Theorem}[]
\newtheorem{corollary}{Corollary}[]
\theoremstyle{definition}
\newtheorem{definition}{Definition}[]
\newtheorem{example}{Example}[]
\theoremstyle{plain}

\newcommand{\aed}{\operatorname{ed_{apr}}}
\newcommand{\ang}{\operatorname{ang}}
\newcommand{\codis}{\operatorname{codis}}
\newcommand{\conv}{\operatorname{conv}}
\newcommand{\diam}{\operatorname{diam}}
\newcommand{\newdiam}{\operatorname{D}}
\newcommand{\dis}{\operatorname{dis}}
\newcommand{\dist}{\operatorname{d}}
\newcommand{\ed}{\operatorname{ed}}
\newcommand{\GH}{\operatorname{\mathcal{G\!H}}}
\newcommand{\id}{\operatorname{id}}
\newcommand{\Int}{\operatorname{Int}}
\newcommand{\Iso}{\operatorname{Iso}}
\newcommand{\Isom}{\operatorname{Isom}}
\newcommand{\dir}{\operatorname{dir}}
\newcommand{\md}{\operatorname{md}}
\newcommand{\tw}{\operatorname{tw}}
\newcommand{\Vol}{\operatorname{Vol}}
\newcommand{\VR}{\operatorname{VR}}

\renewcommand{\:}{\colon}
\newcommand{\0}{\emptyset}
\renewcommand{\c}{\circ}
\newcommand{\rom}[1]{{\em #1}}
\renewcommand{\(}{\rom(}
\renewcommand{\)}{\rom)}
\newcommand{\sm}{\setminus}
\newcommand{\x}{\times}

\newcommand{\e}{\varepsilon}
\renewcommand{\l}{\lambda}
\newcommand{\s}{\sigma}
\renewcommand{\ss}{\subset}

\newcommand{\E}{\mathbb{E}}
\newcommand{\N}{\mathbb{N}}
\newcommand{\I}{\mathbb{I}}
\newcommand{\R}{\mathbb{R}}
\renewcommand{\H}{\mathbb{H}}
\newcommand{\Z}{\mathbb{Z}}


\newcommand{\cM}{\mathcal{M}}
\newcommand{\cP}{\mathcal{P}}
\newcommand{\cR}{\mathcal{R}}

\newcommand{\bA}{\overline{A}}
\newcommand{\bX}{\overline{X}}
\newcommand{\bY}{\overline{Y}}

\newcommand{\Clouds}{\operatorname{Clouds}}
\newcommand{\St}{\operatorname{St}}
\newcommand{\asdim}{\operatorname{asdim}}
\newcommand{\nets}{\operatorname{Net}}

\renewcommand{\U}{\textbf{U}}


\title{Contractibility of the space of $\varepsilon$-nets in $\R$}
\author{Ivan N. Mikhailov}
\date{}

%%%%%%%%%%%%%%%%%%%%%%%%%%%
\begin{document}
\maketitle

\begin{abstract}
In this note, we show that the space of all $\varepsilon$-nets in the real line~$\R$ with a natural metric, equipped with either Hausdorff or Gromov--Hausdorff distance, is contractible.
\end{abstract}
% \begin{abstract}



%%%%%%%%%%%%%%%%%%%%%%%%%%%%%%

%%%%%%%%%%%%%%%%%%%%%%%%%%%%%%

\section{Introduction}

In~\cite{GromovEng} M.\,Gromov introduced moduli spaces of the class of all metric spaces at finite Gromov--Hausdorff distances from a given metric space. It was mentioned that such moduli spaces are always complete and contractible (\cite{GromovEng}[section $3.11_{\frac{1}{2}_+}$]). In~\cite{BogatyTuzhilin} the authors suggested to work with such moduli spaces (they were called \emph{clouds}) in the sense of NBG set theory to avoid arising set-theoretic issues. While the completeness of each cloud was verified in~\cite{BogatyTuzhilin}, the contractibility of each cloud remains an open question for a number of reasons. The main issue here is that a natural homothety-mapping that takes a metric space $(X, d_X)$ into $(X, \lambda d_X)$ for some $\lambda > 0$ and generates a contraction of a cloud of all bounded metric spaces if $\lambda\to 0$, does not behave so well in case of unbounded metric spaces. Firstly, in~\cite{BogatyTuzhilin} it was shown that there exist metric spaces such that $d_{GH}(X, \lambda X) = \infty$ for some $\lambda > 0$. The simplest one is a geometric progression $X = \{3^n\:n\in\N\}$ with a natural metric, for which $d_{GH}(X, 2X) = \infty$. Secondly, even for clouds that are invariant under multiplication on all positive numbers a homothety-mapping may not be continuous. In~\cite{nglitghclitcotrl} it was shown that $d_{GH}(\Z^n, \lambda \Z^n)\ge \frac{1}{2}$ for all $\lambda > 1$, $n\in\N$.

In this note we continue the investigation of the geometry of the Gromov--Hausdorff class. In~\cite{metrictreesRUS} Hausdorff geodesics were constructed, which are geodesic lines in the Gromov--Hausdorff class joining an arbitrary $\varepsilon$-net in $\R$ with $\R$. We show that these natural geodesic lines allow one to construct a contraction of the space of all $\varepsilon$-nets in $\R$, equipped with the Hausdorff metric. We also establish ultrametric inequalities~\ref{thm: UltraNetsHausdorff}, \ref{thm: UltraNetsGromovHausdorff}, which allow us to think of this space as a contractible cone with the vertex in the space $\R$, with the <<Euclidean>> angle at the vertex not exceeding~$\frac{\pi}{2}$, which resembles the situation in the cloud of bounded metric spaces (see Theorem~\ref{thm: boundedcloud}).

% We focus on the other open problem of constructing geodesics in the Gromov--Hausdorff class. In~\cite{Vihrov} a special class of metric spaces in a so-called general position was constructed, such that it is dense in the Gromov--Hausdorff class, and every two metric spaces from these class can be connected with a linear geodesic. However, it is still not known whether every two metric spaces at finite Gromov--Hausdorff distance from each other can be joint with a geodesic in the Gromov--Hausdorff class. Some new examples of geodesics lying in the cloud of the real line appeared in~\cite{nglitghclitcotrl}, and~\cite{IMT} (we review both these constructions in Section~\ref{section: clouds}). In this note, we construct the new geodesic line, for which the real line is an interior point. The existence of such geodesic is impossible in the cloud of bounded metric spaces, because of the ultrametric inequality from the first statement of Theorem~\ref{thm: boundedcloud} below.


\section{Preliminaries}

\emph{A metric space} is an arbitrary pair $(X,\,\dist_X)$, where $X$ is an arbitrary set, $\dist_X\: X\times X\to [0,\,\infty)$ is some metric on it, that is, a nonnegative symmetric, positively definite function that satisfies the triangle inequality.

For convenience, if it is clear in which metric space we are working, we denote the distance between points $x$ and $y$ by $|xy|$. Suppose $X$ is a metric space. By $U_r(a) =\{x\in X\colon |ax|<r\}$, $B_r(a) = \{x\in X\colon |ax|\le r\}$ we denote open and closed balls centered at point~$a$ of radius~$r$ in~$X$. For an arbitrary subset $A\subset X$ of a metric space~$X$, let $U_r(A) = \cup_{a\in A} U_r(a)$ be an open $r$-neighborhood of~$A$. For non-empty subsets $A \ss X$, $B \ss X$ we put $\dist(A,\,B)=\inf\bigl\{|ab|:\,a\in A,\,b\in B\bigl\}$.

\subsection{Hausdorff and Gromov--Hausdorff distances}

\begin{definition}
\label{def: Hausdorff}
Let $A$ and $B$ be non-empty subsets of a metric space.
\emph{The Hausdorff distance} between $A$ and~$B$ is the quantity $$d_H(A,\,B) = \inf\bigl\{r > 0\colon A\subset U_r(B),\,B\subset U_r(A)\bigr\} = \max\left(\sup_{a\in A}\dist(a,\,B),\,\sup_{b\in B}\dist(b,\,A)\right)\!.$$
\end{definition}
\begin{definition} Let $X$ and $Y$ be metric spaces. The triple $(X', Y', Z)$, consisting of a metric space $Z$ and its two subsets $X'$ and $Y'$, isometric to $X$ and~$Y$ respectively, is called \emph{a realization of the pair} $(X, Y)$.
\end{definition}

\begin{definition} \emph{The Gromov-Hausdorff distance} $d_{GH} (X, Y)$ between $X$ and~$Y$ is the exact lower bound of the numbers $r\ge 0$ for which there exists a realization $(X', Y', Z)$ of the pair $(X, Y)$ such that $d_H(X',\,Y') \le r$. 
\end{definition}


Now let $X,\,Y$ be non-empty sets.  

\begin{definition} Each $\sigma\subset X\times Y$ is called a \textit{relation} between $X$ and~$Y$.
\end{definition}

By $\mathcal{P}_0(X,\,Y)$ we denote the set of all non-empty relations between $X$ and~$Y$.

We put $$\pi_X\colon X\times Y\rightarrow X,\;\pi_X(x,\,y) = x,$$ $$\pi_Y\colon X\times Y\rightarrow Y,\;\pi_Y(x,\,y) = y.$$ 

\begin{definition} A relation $R\subset X\times Y$ is called a \textit{correspondence}, if restrictions $\pi_X|_R$ and $\pi_Y|_R$ are surjective.
\end{definition}

Let $\mathcal{R}(X,\,Y)$ be the set of all correspondences between $X$ and~$Y$.

\begin{definition} Let $X,\,Y$ be metric spaces, $\sigma \in \mathcal{P}_0(X,\,Y)$. The \textit{distortion} of $\sigma$ is the quantity $$\dis \sigma = \sup\Bigl\{\bigl||xx'|-|yy'|\bigr|\colon(x,\,y),\,(x',\,y')\in\sigma\Bigr\}.$$
\end{definition}

\begin{proposition}[\cite{BBI}]  \label{proposition: distGHformula}
For arbitrary metric spaces $X$ and~$Y$, the following equality holds $$2d_{GH}(X,\,Y) = \inf\bigl\{\dis\,R\colon R\in\mathcal{R}(X,\,Y)\bigr\}.$$
\end{proposition}

For a metric space~$(X, d_X)$, we denote $\nets_H(X)$ ($\nets_{GH}(X)$) the space of all $\varepsilon$-nets in $X$ equipped with the Hausdorff (Gromov--Hausdorff) distance.

\subsection{Clouds}

Let $\mathcal{VGH}$ denote the class of all nonempty metric spaces endowed with the Gromov--Hausdorff distance.  

\begin{theorem}[\cite{BBI}]
The Gromov--Hausdorff distance is a generalized pseudometric on $\mathcal{VGH}$ that vanishes on every pair of isometric spaces. Namely, the Gromov--Hausdorff distance is symmetric, satisfies the triangle inequality, but in general may be infinite or zero.
\end{theorem}

The class $\mathcal{GH}_0$ is obtained from $\mathcal{VGH}$ by factoring out zero distances, i.e., by the equivalence relation: $X\sim_0 Y$ if and only if $d_{GH}(X,\,Y) = 0$.

\begin{definition}
Consider the equivalence relation $\sim_1$ on $\mathcal{GH}_0$: $X\sim_1 Y$ if and only if $d_{GH}(X,\,Y) < \infty$. The corresponding equivalence classes are called \emph{clouds}.
\end{definition}

For an arbitrary metric space $X$, the cloud defined by it will be denoted by $[X]$. Let $\Delta_1$ denote the metric space consisting of a single point. Thus, $[\Delta_1]$ is the cloud consisting of the classes of all bounded spaces at zero distance from each other.

Suppose that for some pair of metric spaces $A$ and $A'$ we have $d_{GH}(A, A') = 0$. Then for any metric space $B$, the equality $d_{GH}(A, B) = d_{GH}(A', B)$ holds. From this simple observation it follows that any result concerning the Gromov--Hausdorff distance $d_{GH}(A, B)$ remains valid when $A$ is replaced by $A'$ with $d_{GH}(A, A') = 0$. Thus, instead of interpreting the notation $A\in [X]$ directly by definition, assuming that $A$ is an equivalence class of spaces at zero Gromov--Hausdorff distance from $A$, we will, without loss of generality, consider $A$ as a concrete representative of this equivalence class. For example, the notation $X\in[\Delta_1]$ throughout the paper can be read as ``$X$ is a bounded metric space''.

The following theorem on the structure of the cloud $[\Delta_1]$ is integral for this work.

\begin{theorem}[\cite{BBI}] \label{thm: boundedcloud}
Let $X$ and $Y$ be arbitrary bounded metric spaces. Then
\begin{itemize}
\item The following inequalities hold:
$$\frac{1}{2}\bigl|\diam X - \diam Y\bigr|\le d_{GH}(X, Y)\le \max\bigl\{d_{GH}(X, \Delta_1), d_{GH}(Y, \Delta_1)\bigr\} = \frac{1}{2}\max\bigl\{\diam X, \diam Y\bigr\}.$$
\item The map $\Phi\:[\Delta_1]\times \R_{\ge 0}\to [\Delta_1]$, $\Phi(X, \lambda) = \lambda X$ is continuous and yields a contraction of the cloud $[\Delta_1]$ as $\lambda \to 0$.
\item The curve $\lambda X$, $\lambda\in[0, +\infty)$ is a geodesic in the cloud $[\Delta_1]$.
\end{itemize}   
\end{theorem}

\subsection{Auxiliary results}

Finally, we will need two more theorems. The first one allows the construction of Hausdorff geodesics (see \cite{metrictreesRUS} for details):

\begin{theorem}[\cite{metrictreesRUS}, Theorem 9.2]\label{thm:HausdEqGromovHausdR}
Let $X$ be a subset of the real line $\R$, then $d_{GH}(X,\R)=d_H(X,\R)$.
\end{theorem}

And the second one, which shows that a natural homothety-mapping is not a contraction of the cloud $[\R]$:

\begin{theorem}[\cite{nglitghclitcotrl}]\label{thm: homothety-fails}
Consider the integer lattice $\Z^n$, equipped with the Euclidean metric. Then, for all $\lambda > 1$, $n\in\N$, the inequality holds $d_{GH}(\Z^n, \lambda \Z^n)\ge \frac{1}{2}$.   
\end{theorem}

\section{Main results}

\subsection{Hausdorff distance}

\begin{proposition}\label{thm: UltraNetsHausdorff}
For arbitrary $A, B\ss X$, we have $d_H(A, B)\le \max\big\{d_H(A, X), d_H(B, X)\big\}$.
\end{proposition}

\begin{proof}
By Definition~\ref{def: Hausdorff}, we have
\begin{align*}
    &d_H(A, B) = \max\big\{\sup_{a\in A}\dist(a, B), \sup_{b\in B}\dist(b,A)\big\},\\
    &d_H(A, X) = \max\big\{\sup_{a\in A}\dist(a, X), \sup_{x\in X}\dist(x, A)\big\} = \sup_{x\in X}\dist(x, A),\\
    &d_H(B, X) = \max\big\{\sup_{b\in B}\dist(b,X), \sup_{x\in X}\dist(x, B)\big\} = \sup_{x\in X}\dist(x, B).
\end{align*}
Since $\sup_{a\in A}\dist(a, B)\le \sup_{x\in X}\dist(x, B)$ and $\sup_{b\in B}\dist(b, A)\le \sup_{x\in X}\dist(x, A)$, we obtain the desired inequality.
\end{proof}

\begin{example}
Consider the spaces $\lambda_n\Z$ for some sequence $\{\lambda_n\}_{n\in\N}$, tending to~$1$ with $n\rightarrow\infty$. Then spaces $\lambda_N\Z$ do not converge to $\Z$ with respect to the Hausdorff distance. 

Indeed, otherwise it would imply that $\lambda_n\Z\rightarrow_{GH}\lambda\Z$ which contradicts Theorem~\ref{thm: homothety-fails}.
\end{example}

\begin{proposition}
For a finite-dimansional vector space $(V, \|\cdot\|)$, the space $\nets_H(V)$ is contractible. 
\end{proposition}

\begin{proof}
Define $\Phi\: \nets(\R)\times [0, 1]\to \nets(\R)$ as follows: $$\Phi(X, \lambda) = \begin{cases}B_{\lambda/(1-\lambda)}(X) & \lambda\in[0, 1),\\ V &\lambda = 1.\end{cases}$$ 
It suffices to check the continuity of~$\Phi$. Denote $f(x) = \frac{x}{1-x}$.

1) Take arbitrary $\{\lambda_n\}_{n\in\N}$ such that $\displaystyle \lim_{n\to\infty}\lambda_n = \lambda$, where $\lambda\in[0, 1)$. Then $\displaystyle \lim_{n\to\infty}f(\lambda_n) = f(\lambda)$. Note that $d_H\bigl(B_{f(\lambda_n)}(X), B_{f(\lambda)}(X)\bigr)\le |\lambda_n-\lambda|.$ It follows that $\displaystyle \lim_{n\to\infty} B_{f(\lambda_n)}(X) = B_{f(\lambda)}(X)$.

Since $X$ is an $\varepsilon$-net in $V$, there exists $t$ such that for all $t' > t$, we have $B_{t'}(X) = V$. Hence, if $\displaystyle\lim_{n\to\infty}\lambda_n = 1$, then $\displaystyle \lim_{n\to\infty} B_{f(\lambda_n)}(X) = V$.

2) Take arbitrary $\{X_n\}_{n\in\N}$ such that $\displaystyle \lim_{n\to\infty}X_n = X$. Note that $d_H\bigl(B_{f(\lambda)}(X_n), B_{f(\lambda)}(X)\bigr)\le d_H\bigl(X_n, X)$. Indeed, if $p'\in B_{f(\lambda)}(X_n)$, there exist $p\in X_n$ and $q\in X$ such that $\|p-p'\|\le f(\lambda)$, $\|p-q\|\le d_H(X_n, X)$. Thus, $\|p' - (q + p' - p)\| = \|p-q\|\le d_H(X_n, X)$, and $q + p' - p\in B_{f(\lambda)}(X)$, since $\|p-p'\|\le f(\lambda)$. Therefore, since $\displaystyle \lim_{n\to\infty} d_H(X_n, X) = 0$, we obtain $\displaystyle \lim_{n\to\infty} d_H\bigl(B_{f(\lambda)}(X_n), B_{f(\lambda)}(X)\bigr) = 0$, which finishes the proof.
\end{proof}

\subsection{Gromov--Hausdorff distance}

\begin{theorem}\label{thm: UltraNetsGromovHausdorff}
For arbitrary $A,B\ss \R$, we have $d_{GH}(A, B)\le \max\big\{d_{GH}(A, \R), d_{GH}(B, \R)\big\}$.
\end{theorem}
\begin{proof}
Applying theorems~\ref{thm: UltraNetsHausdorff} and~\ref{thm:HausdEqGromovHausdR}, we obtain
$$d_{GH}(A, B)\le d_H(A, B)\le \max\big\{d_H(A, \R), d_H(B, \R)\big\} = \max\big\{d_{GH}(A, \R), d_{GH}(B, \R)\big\}.$$
\end{proof}

\begin{lemma}\label{lemma: transfer_separated}
Suppose $R\in\mathcal{R}(A, B)$, for some subsets $A,B\R$, $\dis R = c$, and $A$ is $t$-separated, for some $t > 2c$. For each $p\in A$, we choose $p'\in R(p)$ arbitrarily. Then, for every three distinct points $p, q, r\in \R$, if $p$ lies between $q$ and~$r$, then $p'$ lies between $q'$ and~$r'$.
\end{lemma}

\begin{proof}
Suppose the desired statement fails for $p,q,r\in \R$. Then, on the one hand, 
\begin{align*}
|q'r'|\ge |qr| - c = |pq| + |pr| - c.
\end{align*}
On the other hand, if $p'$ does not lie between $q'$ and $r'$, then
\begin{align*}
|q'r'| = \bigl||p'q'| - |p'r'|\bigr|\le\max\bigl\{|p'q'| , |p'r'|\bigr\}\le \max\{|pq|, |pr|\} + c \le .
\end{align*}
Thus, we obtain $$2c\ge |pq| + |qr| - \max\bigl\{|pq|, |qr|\bigr\}\ge t, $$ that is a contradiction.
\end{proof}

\begin{theorem}
The space $\nets_{GH}(\R)$ is contractible.
\end{theorem}

\begin{proof}
Define $\Phi\: \nets(\R)\times [0, 1]\to \nets(\R)$ as follows: $$\Phi(X, \lambda) = \begin{cases}B_{\lambda/(1-\lambda)}(X) & \lambda\in[0, 1),\\ \R &\lambda = 1.\end{cases}$$ 
% In other words, $\Phi(X, \lambda) = \overline{\cup_{x\in X}[x-\lambda, x+\lambda]}$. 

\begin{lemma}\label{lemma:contraction}
A mapping $\Phi$ is continuous. 
\end{lemma}

\begin{proof}
It suffices to show that $\Phi$ is continuous with respect to each of two variables $X$ and $\lambda$. Once again, denot $f(x) = \frac{x}{1-x}$.

1) Take arbitrary $\{\lambda_n\}_{n\in\N}$ such that $\displaystyle \lim_{n\to\infty}\lambda_n = \lambda$, where $\lambda\in[0, 1)$. Then $\displaystyle \lim_{n\to\infty}f(\lambda_n) = f(\lambda)$. Consider a natural correspondence $R$ between $B_{f(\lambda_n)}(X)$ and $B_{f(\lambda)}(X)$ which is the union of natural correspondences between segments $[x- f(\lambda_n), x+f(\lambda_n)]$ and $[x-f(\lambda), x + f(\lambda)]$. Note that 
\begin{align*}
\dis R\le 2|f(\lambda_n) - f(\lambda)| = 2\Big|\frac{\lambda_n - \lambda}{(1-\lambda_n)(1-\lambda)}\Big|. 
\end{align*}
 
Thus, $$\lim_{n\to\infty}d_{GH}\Bigl(\Phi(X,\lambda_n), \Phi(X, \lambda)\Bigr) = 0.$$

It remains to show that, for a sequence $\lambda_n$ tending to~$1$, we have $\displaystyle\lim_{n\to\infty} d_{GH}(\Phi(\lambda_n, X), \R) = 0$. Since $X$ is an $\varepsilon$-net in $\R$, we have $d_H(X, \R)<\infty$. Thus, there exists such $t$ that, for any $t' > t$, we have $B_{t'}(X) = \R$, and, hence, the desired statement is true.

2) Take arbitrary $\{X_n\}_{n\in\N}$ such that $\displaystyle \lim_{n\to\infty}X_n = X$. Let us show that $\displaystyle \lim_{n\to\infty}\Phi(X_n, \lambda) = \Phi(X,\lambda)$.

Choose arbitrary $R_n\in\mathcal{R}(X, X_n)$. By passing to a subsequence, we may assume that $\dis R_n \le \frac{2}{n+100}\lambda$. In particular, $d_{GH}(X_n, X)\le \frac{1}{n}$.

Fix some $100$-separated infinite sequence of points $X' = \{p_n\}_{n\in\Z}\ss X$ such that $p_n < p_m \Longleftrightarrow n < m$. Choose $p_k'\in R_n(p_k)$. By Lemma~\ref{lemma: transfer_separated}, the order of points $p_k'$ on the real line is the same (or inverted) as the order of points $p_k$ on the real line. In the remaining proof we assume that these orders coincide with the usual order ($<$) on the real line.

\begin{lemma}\label{lemma: order_change} Suppose $(a, a')\in R_n$, $(b, b')\in R_n$ satisfy inequalities $a < b$, $a' > b'$. Then $b-a \le 2\dis R_n$. 
\end{lemma}

\begin{proof}
Choose $p_k$ such that $p_k > b+100$. Then, by definition of distortion, $\big||p_na|-|p_n'a'|\big|\le \dis R_n$, $\big||p_nb|-|p_n'b'|\big|\le \dis R_n$. Hence, $$0 > b' - a' = |p'-b'| - |p'-a'|\ge |p - b| - |p-a| - 2\dis R_n = b-a - 2\dis R_n \Longleftrightarrow b-a < 2\dis R_n.$$
\end{proof}

Now we extend $R_n$ to a correspondence $R_n'$ between $B_{f(\lambda)}(X)$ and $B_{f(\lambda)}(X_n)$ such that $\dis R_n'\le 5\dis R_n$. 

Take arbitrary $a\in B_{f(\lambda)}(x)$ for some $x\in X$. Unless $a\in\pi_X(R_n)$, we take arbitrary $x'\in R_n(x)$ and add $(a, x' + a - x)$ to $R_n$. 

Let us check the stated inequality. 

We consider three cases:

1) $(a, x' + a-x)$ and $(b, y'+b -y)$, where $a\in B_{f(\lambda)}(x)$, $b\in B_{f(\lambda)}(y)$, $(x, x'), (y, y')\in R_n$.

Here we consider two subcases. 

Wlog, we assume that $x \le y$. 
% (in case $x = y$ everything is clear?). 

\textbf{Case 1.1}. $x' \le y'$. Then 
\begin{multline*}
\big||a-b| - |x'+a-x - y'-b + y|\big| = \big||x + a-x - y - b + y| - |x'+a-x - y'-b + y|\big| \le\\\le \big||x-y| - |x'-y'|\big|\le \dis R_n.
\end{multline*}

\textbf{Case 1.2}. $x' \ge y'$. By Lemma~\ref{lemma: order_change}, we know that $y-x<2\dis R_n < \frac{\lambda}{4}$. Thus,
\begin{align*}
\big||a-b| - |x'+a-x - y'-b + y|\big| \le \big||a-b| - |a-b|\big| + |x-y|+|x'-y'|\le 2|x-y| + \dis R_n\le  5\dis R_n.
\end{align*}


2) $(a, x' + a-x)$ and $(r, r')$, where $a\in B_{f(\lambda)}(x)$, $(x, x'), (r, r')\in R_n$.

This case can be considered similarly to the first one.

3) $(r, r')$ and $(s, s')$ where $(r, r'), (s, s')\in R_n$. 

In this case $\big||r-s|-|r'-s'|\bigr|\le \dis R_n$, by definition of distortion.
\end{proof}


From Lemma~\ref{lemma:contraction}, it follows that $\Phi$ is the desired retraction of $\nets_{GH}(\R)$ onto $\R$.
\end{proof}


%%%%%%%%%%%%%%%%%%%%%%%%%%%%%%%%%%%%
\begin{thebibliography}{2}
% \bibitem{asdim} G.\,Bell, A.\,Dranishnikov, \emph{Asymptotic Dimension}, Topology Appl. 155 (2008),
% 1265–1296.
\bibitem{BogatyTuzhilin} S.\,A.\,Bogaty, A.\,A.\,Tuzhilin, \emph{Gromov--Hausdorff class: its completeness and cloud geometry}, arXiv:2110.06101, 2021.
\bibitem{BBI} D.\,Burago, Yu.\,Burago, S.\,Ivanov, \emph{A Course in Metric Geometry}, Graduate Studies in Mathematics 33, AMS, 2001.
\bibitem{GromovEng} M.\,Gromov, \emph{Metric structures for Riemannian and non-Riemannian spaces}, Birkh\"user
(1999). 
\bibitem{metrictreesRUS} A.\,O.~Ivanov, I.\,N.~Mikhailov, A.\,A.~Tuzhilin, \emph{Gromov--Hausdorff geometry of metric trees}, arXiv:2412.18888, 2024.
\bibitem{nglitghclitcotrl} Mikhailov.\,I., \emph{New geodesic lines in the Gromov--Hausdorff class lying in the cloud of the real line}, Chebyshevskii Sbornik. 2025;26(2):176-185. (In Russ.); arXiv:2504.14629, 2025.
\bibitem{TuzhilinHistory} A.\,A.\,Tuzhilin, \emph{Who invented the Gromov-Hausdorff Distance?}, 2016, ArXiv e-prints, arXiv:1612.00728.




\end{thebibliography}


\end{document} 
